\section{Geometrische Optik}

\textbf{Grundlagen:}\\
\hspace{5mm} - geradlinige Ausbreitung von Licht (Lichtstrahlen) \\
\hspace{5mm} - Vernachlässigung von Wellenphänomenen (Beugung und Interferenz)

\subsection{Spiegel}

\textbf{Grundbegriffe:}
\begin{itemize}
	\item \textbf{virtuelles Bild:} Im Bildpunkt treffen sich nur die Verlängerungen der Strahlen
	\item \textbf{reelles Bild:} Im Bildpunkt treffen sich tatsächlich Lichtstrahlen
\end{itemize}

\textbf{Variablen:}
\begin{itemize}
	\item $G$: Gegenstandshöhe
	\item $g$: Gegenstandsweite
	\item $B$: Bildhöhe
	\item $b$: Bildweite
\end{itemize}

\textbf{Ebener Spiegel:} $G = B$, $g = b$ \\

\textbf{Sphärischer Spiegel:}
\begin{formulabox}
    \[
        \frac{1}{g} + \frac{1}{b} = \frac{2}{r}
    \]
\end{formulabox}
für $g \rightarrow \infty$ folgt: $b = \frac{r}{2} = f$: Brennweite des Spiegels \\
Parallel einfallende Strahlen fokussieren in der Brennebene \\
Achsenparallele Strahlen fokussieren im Brennpunkt \\[1ex]

\textbf{Bildkonstruktion bei sphärischen Spiegeln} \\
1. Der achsenparallele Strahl wird in den Brennpunkt reflektiert \\
2. Der Brennpunktstrahl verläuft durch den Brennpunkt und wird achsenparallel \\ \hspace{2.3mm} reflektiert \\
3. Der Mittelpunktstrahl verläuft durch den Krümmungsmittelpunkt, trifft den Spiegel und \\ \hspace{2.3mm} wird in sich selbst zurück reflektiert \\[1ex]

\textbf{Abbildungsmassstab:}
\begin{formulabox}
    \[
        V = \frac{B}{G} = -\frac{b}{g}
    \]
\end{formulabox}
$V > 0$ das Bild steht aufrecht\\
$V < 0$ das Bild ist invertiert \\[1ex]

\textbf{Vorzeichenregeln:} \\
1. $g$ ist positiv, wenn der Gegenstand auf der Seite des Spiegels ist, von der aus Licht \\ \hspace{2.3mm} auf den Spiegel trifft \\
2. $b$ ist positiv, wenn sich das Bild auf der Seite des Spiegels befindet, in die das Licht \\ \hspace{2.3mm} vom Spiegel reflektiert wird \\
3. $r$ (und daher auch $f$) ist positiv, wenn der Spiegel konkav ist. Dann liegt der \\ \hspace{2.3mm} Brennpunkt auf der Seite des Spiegels von der Licht einfällt 

\subsection{Linsen}
Ann.: Sphärische Grenzflächen zwischen zwei Medien mit Brechzahlen $n_1$ und $n_2$ \\[1ex]
\textbf{Brechung an einer Grenzfläche}
\begin{formulabox}
    \[
        \frac{n_1}{g} + \frac{n_2}{b} = \frac{n_2 - n_1}{r}
    \]
\end{formulabox}

\textbf{Vorzeichenkonvention für die Brechung:} \\
1. $g$ ist positiv, für Gegenstände auf der Einfallsseite der brechenden Fläche \\
2. $b$ ist positiv, für Bilder auf der Transmissionsseite der brechenden Fläche \\
3. $r$ ist positiv, wenn der Krümmungsmittelpunkt auf der Transmissionsseite der brechenden Fläche liegt

\subsection{Dünne Linsen}
\textbf{Für die Brechung an der ersten Grenzfläche gilt:}
\[
    \frac{n_{Luft}}{g} + \frac{n_{Linse}}{b_1} = \frac{n_{Linse} - n_{Luft}}{r_1}
\]
mit $\Rightarrow g_2 = -b_1$ folgt für die \textbf{Brechung an der zweiten Grenzfläche:}
\[
    \frac{n_{Linse}}{-b_1} + \frac{n_{Luft}}{b} = \frac{n_{Luft} - n_{Linse}}{r_2}
\]
\begin{formulabox}
    \[
        \frac{1}{g} + \frac{1}{b} = \left( \frac{n_{Linse}}{n_{Luft}} -1 \right) \left( \frac{1}{r_1} - \frac{1}{r_2} \right)
    \]
\end{formulabox}

\textbf{Reziproke Brennweite einer dünnen Linse:}
\begin{formulabox}
    \[
        \frac{1}{f} = \left( \frac{n_{Linse}}{n_{Luft}} -1 \right) \left( \frac{1}{r_1} - \frac{1}{r_2} \right)
    \]
\end{formulabox}

\textbf{Allgemeine Abbildungsgleichung:}
\begin{formulabox}
    \[
        \frac{1}{g} + \frac{1}{b} = \frac{1}{f}
    \]
\end{formulabox}

\textbf{Bildkonstruktion mit Linsen} \\
1. Der achsenparallele Strahl wird so gebrochen, dass er durch den zweiten \\ \hspace{2.3mm} Brennpunkt der Linse verläuft \\
2. Der Mittelpunktstrahl verläuft durch den Linsenmittelpunkt und wird nicht \\ \hspace{2.3mm} gebrochen (nur bei dünner Linse) \\
3. Der Brennpunktstrahl verläuft durch den ersten Brennpunkt und verlässt die \\ \hspace{2.3mm} Linse achsenparallel \\[1ex]

\textbf{Abbildungen mit mehreren Linsen} \\
Für zwei dünne Linsen auf derselben optischen Achse gilt:
\begin{formulabox}
    \[
        \frac{1}{f} = \frac{1}{f_1} + \frac{1}{f_2}
    \]
\end{formulabox}
Sie wirken wie eine dünne Linse mit der effektiven Brennweite $f$
